%%%% CAPÍTULO 1 - INTRODUÇÃO
%% Problema, motivação, objetivos e organização do relatório.

\chapter{Introdução}\label{cap:introducao}

A disseminação de notícias falsas (\textit{fake news}) em redes sociais é um fenômeno de difusão de informação estudado por diferentes áreas, da sociologia à física estatística~\cite{Vosoughi2018}. Modelos computacionais permitem simular essa propagação em grandes populações, mas o custo de simular milhões de indivíduos ao longo de muitas gerações torna necessário o uso de computação paralela e distribuída.

Este trabalho parte de uma simulação sequencial, escrita em Python e fornecida na disciplina, em que a população é representada por uma matriz bidimensional e cada indivíduo assume um de três estados (ignorante, espalhador ou inativo), com propagação local pela vizinhança de Moore. O objetivo é transformar essa aplicação em versões paralela e distribuída, preservando o comportamento lógico da versão original, e avaliar o desempenho obtido.

\todo{Completar a motivação com 1--2 parágrafos e referências.}

\section{Objetivos}

\subsection{Objetivo geral}

Desenvolver e avaliar versões paralela e distribuída de uma simulação de propagação de \textit{fake news} baseada em autômato celular, comparando-as com a versão sequencial em termos de \textit{speedup}, eficiência e limitações.

\subsection{Objetivos específicos}

\begin{itemize}
    \item Compreender e documentar o modelo computacional da versão sequencial;
    \item Implementar uma versão paralela com divisão explícita da matriz entre \textit{threads}/processos, sem condições de corrida;
    \item Implementar uma versão distribuída com \textit{sockets}, com sincronização das fronteiras da matriz entre processos;
    \item Verificar que as três versões produzem resultados idênticos para os mesmos parâmetros;
    \item Propor e avaliar melhorias no modelo e na implementação;
    \item Medir e analisar desempenho variando tamanho da matriz, gerações, percentual inicial de espalhadores, número de \textit{threads} e de processos/\textit{hosts}.
\end{itemize}

\section{Organização do relatório}

%% NBR 14724:2024: o trabalho se divide em seções (NBR 6024), não em "capítulos".
%% Por isso usar "Seção~\ref{...}" em vez de \autoref (que imprime "Capítulo").
A Seção~\ref{cap:fundamentacao} apresenta a fundamentação teórica. A Seção~\ref{cap:sequencial} descreve o modelo e a versão sequencial. As seções~\ref{cap:paralela} e~\ref{cap:distribuida} apresentam as versões paralela e distribuída. A Seção~\ref{cap:melhorias} descreve as melhorias implementadas, a Seção~\ref{cap:experimentos} apresenta os experimentos e resultados, e a Seção~\ref{cap:conclusoes} traz as conclusões.
